\section{Implémentation du réseau Wi-Fi HaLow}

La norme IEEE 802.11 intègre intrinsèquement des mécanismes de fiabilité, tels que les accusés 
de réception (ACK) et les retransmissions automatiques, destinés à garantir l'intégrité des données. 
Cependant, dans le cadre d'un flux vidéo temps réel pour le FPV, la perte d'une trame est visuellement 
préférable à l'attente d'une retransmission, qui induit une latence beaucoup trop élevée. L'architecture 
réseau a donc été configurée pour s'affranchir de ces contraintes.

\subsection{Stratégie de diffusion (UDP Broadcast et No-ACK)}
Pour contourner l'attente des acquittements matériels de la couche MAC\index{Couche MAC}, l'application 
embarquée sur le STM32 utilise le protocole de transport UDP sur une adresse IP de diffusion \textit{Broadcast}.

Dans le code C exploitant la pile LwIP, l'adresse de destination est configurée sur l'adresse de diffusion 
locale du sous-réseau :
\begin{lstlisting}[language=C]
IP4_ADDR(ip_2_ip4(&dest_ip), 192, 168, 12, 255);
\end{lstlisting}

En ciblant l'adresse \texttt{.255}, la couche réseau LwIP force l'en-tête Wi-Fi à utiliser l'adresse MAC 
de destination \texttt{FF:FF:FF:FF:FF:FF}. Selon la norme 802.11, les trames diffusées à cette adresse 
ne déclenchent aucun acquittement (ACK) de la part des récepteurs. L'émetteur expédie ainsi les paquets 
en flux tendu, ce qui annule la latence induite par les retransmissions au prix de potentielles pertes 
de paquets.

De plus, la taille de la charge utile (\textit{Payload}) UDP a été fixée à 1400 octets :
\begin{lstlisting}[language=C]
struct pbuf *p = pbuf_alloc(PBUF_TRANSPORT, 1400, PBUF_REF);
\end{lstlisting}
Cette valeur approche la taille maximale d'une trame Ethernet standard (MTU), ce qui minimise le ratio 
d'en-têtes (Overhead MAC+IP+UDP) par rapport aux données utiles.

\subsection{Optimisation de la réactivité radio : Désactivation du PSM}
Par défaut, les modules Wi-Fi HaLow Morse Micro intègrent des mécanismes de gestion de l'énergie 
(\textit{Power Saving Mode}) destinés aux capteurs IoT sur batterie. Ce mode place la radio en 
veille entre les transmissions, ce qui induisait dans nos tests une latence systématique de environ 
\textbf{22 ms} par paquet.

Pour un usage FPV, cette économie d'énergie n'est pas accéptable. Le mode veille a été explicitement désactivé 
dans le firmware du STM32 via l'appel suivant :
\begin{lstlisting}[language=C]
mmwlan_set_power_save_mode(MMWLAN_PS_DISABLED);
\end{lstlisting}
Cette modification a permis de stabiliser le temps de réponse réseau (\textit{ping}) entre \textbf{3 et 5 ms}, 
garantissant une gigue minimale.

\subsection{Contournement des contraintes réglementaires (ETSI vs FCC)}
Lors des premiers essais d'établissement de la liaison radio en MCS 1, une latence de base colossale 
d'environ 670 ms a été observée sur de simples paquets de test (Ping). 

L'analyse de ce délai a mis en évidence l'impact restrictif des normes réglementaires européennes (ETSI) 
régissant la bande des 868 MHz. Ces normes imposent un mécanisme de partage des ondes strict appelé 
\textit{Listen Before Talk} (LBT) ainsi que des restrictions de \textit{Duty Cycle} (temps d'émission 
maximal autorisé). Le module radio retenait intentionnellement les paquets pour respecter ces seuils 
légaux.

Afin d'évaluer les performances brutes du changement de MCS sans ces bridages logiciels imposés par la 
région, la configuration de l'émetteur et du récepteur a été basculée sur le domaine réglementaire 
américain (FCC, bande des 915 MHz). L'utilisation du Canal 27 (915.5 MHz) a permis de désactiver 
les mécanismes LBT, ramenant immédiatement la latence réseau de base à environ 27 ms pour des 
paquets de 1450 octets.

\subsection{Le routage et l'isolation IP (OpenWrt)}
Pour éviter les conflits de routage sur la station de réception, le réseau FPV est strictement isolé 
sur le sous-réseau \texttt{192.168.12.x}.

Lors de l'implémentation, un conflit IP majeur bloquait la transmission : le routeur OpenWrt (récepteur) 
possédait l'adresse \texttt{192.168.12.1} simultanément sur son interface Wi-Fi HaLow (\texttt{ahwlan}) 
et sur son interface Ethernet locale (\texttt{lan}). Pour résoudre ce conflit, il a suffi de migrer le 
LAN sur \texttt{192.168.13.1}.

Concernant le transfert des données à travers le routeur, une première approche a été 
testée : les paquets UDP entrants sur l'interface sans-fil étaient interceptés (via \texttt{tcpdump}) 
puis redirigés manuellement vers le port Ethernet à l'aide de l'utilitaire \texttt{netcat} (\texttt{nc}). 
Cependant, ce traitement en espace utilisateur (\textit{User-Space}) sollicitait lourdement le processeur 
du routeur et induisait une latence importante.
\begin{lstlisting}[language=bash]
tcpdump -i wlan0 -U -s 0 -w - "udp port 1337" | nc 192.168.12.10 9999
\end{lstlisting}

Pour optimiser le routage, un pont logiciel (\textit{Bridge}) a finalement été forcé au niveau du système 
d'exploitation du routeur :
\begin{lstlisting}[language=bash]
brctl addif br-ahwlan wlan0
\end{lstlisting}
Ce pont matériel permet au flux \textit{Broadcast} en provenance de l'interface sans-fil de traverser le 
routeur et d'être répliqué directement sur le câble Ethernet(pour autant que la station de réception
connectée soit sur le bon sous réseau) sans être intercepté ni traité par la pile 
TCP/IP ou des utilitaires logiciels, économisant ainsi 10 à 20 ms de délai de traitement additionnel.

\subsection{Validation empirique et mathématique de la bande passante vidéo}
La désactivation de l'adaptation automatique de débit (\textit{Rate Control}) a été effectuée dans le 
SDK via la fonction ATE (\textit{Automated Test Equipment}, fonction utilisée en interne pour forcer 
des configurations matérielles spécifiques lors des phases de tests) :
\begin{lstlisting}[language=C]
mmwlan_ate_override_rate_control(MMWLAN_MCS_2, MMWLAN_BW_2MHZ, MMWLAN_GI_NONE);
\end{lstlisting}

En forçant l'utilisation du \textbf{MCS 2} sur une largeur de bande de \textbf{2 MHz}, la modulation 
employée (QPSK) garantit une excellente robustesse face aux obstacles physiques. Des mesures empiriques 
(par capture \texttt{tcpdump}) ont validé un débit utile réseau stable d'environ \textbf{1.34 Mbps} 
(soit une moyenne de 120 paquets de 1400 octets par seconde).

Il est ensuite nécessaire de confronter cette capacité physique aux exigences du flux vidéo. Le débit brut 
(non compressé) se calcule ainsi :
$$D_{brut} = W \times H \times F \times B_{pp}$$

Afin de contourner la rétention d'image du processeur de signal (ISP) de la caméra, la fréquence 
d'échantillonnage a dû être fixée à 60 FPS. Pour le flux vidéo retenu en résolution 480p 
(soit $720 \times 480$ pixels), capturé à 60 images par seconde et codé sur 24 bits de couleurs :
$$D_{brut\_480p60} = 720 \times 480 \times 60 \times 24 = 497\,664\,000 \text{ bps}$$
Soit environ 497.6 Mbps.

Un encodeur matériel moderne (H.264 / H.265), configuré pour du streaming en temps réel, offre un ratio de 
compression effectif standard avoisinant 1:200. Le débit réseau théorique requis deviendrait alors :
$$D_{cible} = \frac{D_{brut}}{C_{ratio}}$$
$$D_{cible\_480p60} = \frac{497.6 \text{ Mbps}}{200} \approx 2.49 \text{ Mbps}$$

Cette modélisation démontre clairement qu'à 60 images par seconde, même une résolution modeste de 480p avec 
une compression standard exigerait environ 2.49 Mbps, ce qui saturerait immédiatement le canal radio imposé 
par le MCS 2 (limité empiriquement à 1.34 Mbps).

C'est pour cette raison critique que l'encodeur matériel doit être paramétré de manière drastique. Afin de 
forcer ce flux à traverser la liaison radio tout en conservant une marge de sécurité vitale face aux baisses 
imprévisibles de la qualité du lien RF en condition de vol NLOS, le taux de compression final de l'encodeur 
a été bridé de manière stricte (CBR - \textit{Constant Bit Rate}) à \textbf{400 kbps}. 

\subsubsection{Impact de la compression sur la fragmentation des paquets}
Ce compromis de 400 kbps impose à l'encodeur un ratio de compression extrême (environ 1:1240) au détriment 
de la fidélité visuelle, mais garantit un avantage décisif sur la latence matérielle. En lissant ce débit 
sur la fréquence d'échantillonnage, le poids moyen d'une image vidéo compressée se calcule ainsi :
$$Taille_{image} = \frac{400\,000 \text{ bits/s}}{60 \text{ FPS} \times 8 \text{ bits}} \approx 833 \text{ octets}$$

Sachant que la charge utile maximale (MTU) définie pour nos trames UDP et nos transferts SPI est de 
\textbf{1400 octets}, l'écrasante majorité des images générées par l'encodeur tiennent dans 
\textbf{un seul et unique paquet réseau}. Cette caractéristique est fondamentale pour la basse latence : 
elle évite la fragmentation IP et exempte le récepteur de devoir stocker et attendre de multiples paquets 
pour reconstruire et décoder une seule trame vidéo.